% Generated draft: requires mathematical and source review.
\documentclass[11pt]{article}
\usepackage[margin=25mm]{geometry}
\usepackage{fontspec}
\setmainfont{Latin Modern Roman}
\usepackage{amsmath,amssymb,mathtools,unicode-math}
\setmathfont{Latin Modern Math}
\usepackage{enumitem,needspace,xurl,hyperref}
\setlength{\parindent}{0pt}
\setlength{\parskip}{5pt}
\setlength{\emergencystretch}{3em}
\newcommand{\N}{\mathbb{N}}
\newcommand{\Z}{\mathbb{Z}}
\newcommand{\Q}{\mathbb{Q}}
\newcommand{\R}{\mathbb{R}}
\newcommand{\C}{\mathbb{C}}
\newcommand{\F}{\mathbb{F}}
\newcommand{\A}{\mathbb{A}}
\newcommand{\PP}{\mathbb{P}}
\newcommand{\E}{\mathbb{E}}
\newcommand{\eps}{\varepsilon}
\newcommand{\dd}{\,\mathrm{d}}
\begin{document}
\section*{UnsolvedMath ID: 11000034}
\textbf{Source problem code:} catalogue entry
\subsection{Definitions and mathematical statement}
Let $\operatorname{Mod}_g$ be the mapping class group of a closed orientable genus-$g$ surface, and let $S$ be a finite symmetric generating set. For the word metric $|\cdot|_S$, put
\[
 B_{G,S}(r)=\{x\in G:|x|_S\le r\},\qquad
 F_{G,S}(z)=\sum_{r\ge0}|B_{G,S}(r)|z^r.
\]
Say that $G$ has rational growth for $S$ if $F_{G,S}(z)$ is the Taylor series of a rational function in $\Q(z)$ regular at $0$. The sphere-counting convention gives the same rationality question, because its series equals $(1-z)F_{G,S}(z)$.

\textbf{Question 3.13 (rational growth).} Does $\operatorname{Mod}_g$ have rational growth with respect to some finite generating set? Does it have rational growth with respect to every finite generating set? These are distinct quantifiers:
\[
 \exists S\quad F_{\operatorname{Mod}_g,S}(z)\in\Q(z),
 \qquad\text{and}\qquad
 \forall S\quad F_{\operatorname{Mod}_g,S}(z)\in\Q(z).
\]
Ask the corresponding questions for any finitely generated subgroup $H<\operatorname{Mod}_g$, with its own finite generating set and intrinsic word metric, including the Torelli group $\mathcal I_g$ in the genus range where that example is finitely generated. The excerpt supplies no genus restriction.

The source then introduces property-counting series. For a specified property $\mathcal P$ of elements of $G=\operatorname{Mod}_g$, define
\[
 c_{\mathcal P,S}(r)=\#\{x\in B_{G,S}(r):x\text{ has }\mathcal P\},\qquad
 f_{\mathcal P,S}(z)=\sum_{r\ge0}c_{\mathcal P,S}(r)z^r.
\]
Its examples are finite order, membership in a fixed subgroup, and being pseudo-Anosov. The coefficients count elements in balls, not words, spheres, or conjugacy classes. This is context for seeking other recursive patterns; the excerpt does not assert that the series is rational for every arbitrary property.

The source emphasizes that an automatic structure by itself is not known there to imply rational growth, and that a geodesic automatic structure supplies the stronger relevant condition. It refers to Mosher's automatic structure and says that this particular structure is not geodesic. The registry summary instead says that $\operatorname{Mod}_g$ is not automatic for $g\ge2$. Those supplied assertions conflict: the summary cannot be adopted as an explanation of the source question without correction and primary-source checking.
\subsection{Short English statement}
Ask whether mapping class groups, and their finitely generated subgroups, have rational growth for some or every finite generating set. The source also defines ball-counting growth series for properties such as torsion, subgroup membership, and pseudo-Anosov behavior.
\subsection{Sources}
\begin{itemize}
\item Local UnsolvedMath registry, permanent ID 11000034, source code catalogue entry; no individual proposer or proposal year supplied.
\item Farb, edited volume \emph{Problems on Mapping Class Groups and Related Topics}, Question 3.13, PDF page 35; \url{https://www.math.uchicago.edu/~farb/papers/mcgbook.pdf}. The registry records PDF-font-structure extraction and access on 2026-07-29; the inventory distinguishes 329 labeled headings by source page.
\item The source mentions Mosher's automatic structure without a complete citation in this record. The registry summary's automaticity statement conflicts with this source discussion and has no supporting citation supplied.
\end{itemize}
\end{document}
